\chapter{Induced Acceleration Analysis}
\label{ch:induced_acceleration}
\raggedbottom

\begin{laymansbox}{The Question Behind This Chapter}{lay:iaa_question}
A golfer's wrist can accelerate even when the net wrist moment is zero. Conversely, a large wrist moment need not produce a large increase in clubhead speed. The explanation lies in the coupled mechanics of the whole system: mass distribution, posture, current motion, external support and the point whose acceleration we ask about.

\textbf{Induced acceleration analysis} (IAA) computes how a specified force term enters a model's instantaneous acceleration. It connects the force accounting of earlier chapters to motion without mistaking a sequence of segment-speed peaks for a measured sequence of energy transfers. Its useful question is precise: at this state, in this model, how do the declared inputs and other terms combine to produce the chosen output acceleration?

The method has a substantial biomechanics history, including the review by \citep{ZajacGordon1989}, throwing analysis by \citep{Hirashima2008kinetic}, and walking simulations by \citep{Neptune2001}. Golf also has prior acceleration-partition research: \citep{Takagi2021pelvis} analyzed pelvis rotation in skilled golfers. This chapter develops the mechanics and their connection to affine control; it does not claim the first application to golf.
\end{laymansbox}

\begin{infobox}{Normative Attribution Record}
A publishable attribution must identify the model and revision; engine, solver, and revision; generalized coordinates and reference frame; output point; mass matrix; constraint handling; contact model and mode; complete force partition; residual treatment; and numerical tolerance with closure error. Its identifiability contract must state the measurements or assumptions supporting each term. An algebraic term does not by itself identify anatomical source, neural intent, necessity, sufficiency, or intervention effect. Missing fields make a result \textbf{unsupported or unqualified}; an unavailable cross-engine state is reported that way rather than treated as solver parity.

Two counterexamples block unique attribution. Under $q=Tz$, the same dynamics use $M_z=T^\mathsf{T}M_qT$ and $Q_z=T^\mathsf{T}Q_q$: component values change even though $T\ddot z=\ddot q$. Likewise, if $Q=Q_A+Q_B$, then $Q=(Q_A+r)+(Q_B-r)$: a residual reassignment changes both reported terms while preserving their sum. An intervention must separately re-solve constraints and contact and declare what is held fixed.
\end{infobox}

\section{The Mathematical Foundation}
\label{sec:iaa_math}

\subsection{Starting From the Manipulator Equation}
First consider independent generalized coordinates \(q\), with generalized velocity \(v=\dot q\). Our sign convention puts all retained non-inertial terms on the right:
\begin{equation}
 M(q)\ddot q=Q_u+Q_v+Q_g+Q_p+Q_e+Q_r.
 \label{eq:iaa_eom}
\end{equation}
Here \(Q_u\) denotes applied control, \(Q_v\) the velocity-dependent inertial terms, \(Q_g\) gravity, \(Q_p\) passive constitutive forces, \(Q_e\) specified external loads, and \(Q_r\) residual forces used to reconcile model and data. These categories must not overlap. If a conventional equation uses \(M\ddot q+C(q,v)v+g(q)=Bu+Q_e\), then our convention has \(Q_v=-Cv\), \(Q_g=-g\), and \(Q_u=Bu\).

For a regular independent-coordinate model, \(M\) is symmetric positive definite. Solving the coupled system gives
\begin{equation}
 \ddot q=M^{-1}Q,\qquad Q=\sum_k Q_k.
 \label{eq:iaa_forward}
\end{equation}
Implementations should factor and solve with \(M\), rather than explicitly form its inverse. The inverse notation identifies the mathematical operator. A floating base can have an invertible mass matrix even though some coordinates are unactuated; underactuation and singular inertia are different conditions.

\begin{equation}
 \ddot q=\sum_k a_k,\qquad a_k=M^{-1}Q_k.
 \label{eq:iaa_superposition}
\end{equation}
This superposition is exact at the \emph{same} state and parameters. It does not add together trajectories simulated under different inputs. Nor does it authorize dividing a joint torque by that joint's isolated inertia: the full coupled solve is essential.

\subsection{What a Muscle Label Requires}
Net joint moments inferred from inverse dynamics are not individual muscle forces. If a muscle model supplies forces \(F_i\) and generalized moment-arm columns \(r_i(q)\), its actuator terms are \(Q_i=r_iF_i\). Many muscle-force combinations can generate the same net moment, including antagonistic co-contraction. Optimization, activation dynamics, EMG constraints and other measurements can narrow that ambiguity, but their assumptions remain part of the result.

The choice of boundary matters equally. A grip wrench is external to a club-only model and internal to a combined golfer--club model. Counting it as both a specified load and a separately solved internal reaction duplicates the same interaction. The reader must be able to trace every force to one place in the accounting.

\subsection{Constraints Add a Baseline}
Suppose a smooth, fixed contact mode imposes independent acceleration constraints \(A\ddot q=b\). For a stationary holonomic constraint \(\phi(q)=0\), \(A=\partial\phi/\partial q\) and \(b=-\dot A v\). The state must already satisfy the position and velocity constraints. With the reaction convention \(A^{\mathsf T}\lambda\), solve
\[
 \begin{aligned}
 M\ddot q&=Q+A^{\mathsf T}\lambda,\\
 A\ddot q&=b.
 \end{aligned}
\]
Substitution yields the Schur complement \(W=AM^{-1}A^{\mathsf T}\) and
\[
 \begin{aligned}
 \lambda&=W^{-1}(b-AM^{-1}Q),\\
 P&=M^{-1}-M^{-1}A^{\mathsf T}W^{-1}AM^{-1},\\
 a_c&=M^{-1}A^{\mathsf T}W^{-1}b,\\
 \ddot q&=PQ+a_c.
 \end{aligned}
\]
These formulas assume positive-definite \(M\) and full row rank of \(A\). Redundant constraints require a consistent rank treatment, not an unexplained inverse of a singular matrix. Quaternion coordinates also require their proper velocity map; the simplifying convention \(v=\dot q\) is not universal.

The force-to-acceleration map is \textbf{affine}, because \(a_c\) need not vanish. The correct partition is \(\ddot q=\sum_k PQ_k+a_c\): count the baseline once. Each source increment satisfies \(A(PQ_k)=0\), whereas \(Aa_c=b\). Solving a full constrained problem for every isolated force and summing those solutions would count the curvature baseline repeatedly. Although \(P\) is symmetric positive semidefinite, it is not generally an idempotent Euclidean projection; it has inverse-inertia units.

For example, a 2 kg particle at \(q=(1,0)\) m on a fixed 1 m circular guide has velocity \(v=(0,3)\) m/s. Taking \(A=(1,0)\), the acceleration constraint is \(\ddot q_x=-9\) m/s\(^2\). The solution is
\[
 \ddot q=(-9,F_y/2)^{\mathsf T},\qquad \lambda=-18-F_x.
\]
The radial acceleration is present even with zero applied force, while changing \(F_x\) changes the reaction. This is a bilateral guide example. A unilateral contact must also pass its reaction-sign and friction-capacity conditions. If the required contact cannot persist, the active mode must change and this frozen operator no longer describes the realizable response.

\subsection{From Joint Acceleration to Clubhead Acceleration}
Let \(p(q)\) be a point expressed in a fixed inertial frame and \(J=\partial p/\partial q\). Then
\[
 \ddot p=J\ddot q+\dot Jv
         =\sum_k JPQ_k+Ja_c+\dot Jv.
\]
The last term is the kinematic curvature of the output map. Joint angular acceleration can be zero while a rotating clubhead has substantial centripetal acceleration. Thus a joint-coordinate ledger cannot simply be renamed a clubhead-acceleration ledger.

For two force choices at the \emph{same} state and mode, the baseline cancels: \(\Delta\ddot p=JP\Delta Q\). This is often the clearest sensitivity to report. Use a feasible reference load when zero force would violate the contact conditions. For clubhead speed, project the total acceleration along its velocity: \(\dot s=\hat v_p^{\mathsf T}\ddot p\) when \(s=\|\dot p\|>0\). A large acceleration normal to motion bends the path without increasing speed at that instant. At zero speed this unit-tangent formula is undefined.

\subsection{Connection to the Control-Affine Framework}
\label{sec:iaa_affine_connection}
When \(Q=Q_0(q,v)+B(q)u\), the local constrained equations give
\[
 \frac{d}{dt}\begin{bmatrix}q\\v\end{bmatrix}
 =\underbrace{\begin{bmatrix}v\\PQ_0+a_c\end{bmatrix}}_{f(x)}
 +\underbrace{\begin{bmatrix}0\\PB\end{bmatrix}}_{G(x)}u.
\]
IAA and affine control share this frozen-state input map when they describe the same plant and input. The drift includes the kinematic block \(v\), all retained autonomous forces and the constraint baseline. It is not just a velocity-dependent torque. Passive tissue forces and muscle activation or tendon states can enlarge the autonomous system. A generalized-moment input is not interchangeable with neural excitation; whether the enlarged plant is affine in that excitation must be derived from its actuator equations.

Contact switches, impacts and input-dependent mode changes require a hybrid or nonsmooth description. The equation above is a mode-local statement, not a global proof of control affinity through every phase of a golf swing.

\section{Dynamic Coupling and Posture}
\label{sec:dynamic_coupling}

\subsection{Why One Torque Can Accelerate Several Coordinates}
Write \(H=M^{-1}\) in the unconstrained case. For a two-coordinate torque increment,
\begin{equation}
 \delta\ddot q=H\begin{bmatrix}\delta\tau_1\\0\end{bmatrix}
 =\begin{bmatrix}H_{11}\delta\tau_1\\H_{21}\delta\tau_1\end{bmatrix}.
 \label{eq:iaa_coupling}
\end{equation}
Off-diagonal entries allow remote acceleration. They are often nonzero, but physical connection does not guarantee a nonzero response at every coordinate in every posture. The analogous constrained map is \(P\), and the physical point response uses \(JP\).

Consider two uniform planar rods, each of length 1 m and mass 1 kg, pinned at the proximal end. Their center-of-mass inertias are \(1/12\) kg m\(^2\). Let \(q_1\) measure the first rod from vertically downward, and \(q_2\) the second rod relative to the first. The mass matrix, in kg m\(^2\), is
\[
 M(q_2)=\begin{bmatrix}
 5/3+\cos q_2&1/3+\tfrac12\cos q_2\\
 1/3+\tfrac12\cos q_2&1/3
 \end{bmatrix}.
\]
At \(q_2=0\), a unit proximal torque increment produces \((12/7,-30/7)\) rad/s\(^2\); a unit distal torque produces \((-30/7,96/7)\) rad/s\(^2\). The remote gains are equal because \(H\) is symmetric. At \(q_2=\arccos(-2/3)\), however, the matrix is diagonal and the same proximal increment produces \((1,0)\) rad/s\(^2\). Connected rods can have zero instantaneous cross-acceleration. These manufactured states establish mechanics, not an optimal golf posture.

\subsection{Normalized Coupling Is Not Nonreciprocal Mobility}
Challis introduced an induced-acceleration index for joint coupling \citep{Challis2011}. To avoid an ambiguous inverse-element convention, we derive a specifically defined two-coordinate ratio here. Under a force increment only at coordinate 1, the unforced second row gives \(M_{21}\delta\ddot q_1+M_{22}\delta\ddot q_2=0\). Therefore
\begin{equation}
 R_{1\to2}=\left|\frac{\delta\ddot q_2}{\delta\ddot q_1}\right|
 =\left|\frac{H_{21}}{H_{11}}\right|
 =\left|\frac{M_{21}}{M_{22}}\right|.
 \label{eq:iai}
\end{equation}
The denominator is the \emph{mass-matrix entry} \(M_{22}\), not the inverse-matrix entry \(H_{22}\). This ratio excludes the common acceleration baseline. With multiple unforced coordinates indexed by \(\mathcal P\), use the block relation
\[
 \delta\ddot q_{\mathcal P}
 =-M_{\mathcal P\mathcal P}^{-1}M_{\mathcal P1}\delta\ddot q_1.
\]
Independent elementwise divisions generally fail because those passive coordinates couple to one another. Additional constraints require the constrained operator instead of this unrestricted block formula.

For the straight two-rod model, \(R_{1\to2}=2.5\) and \(R_{2\to1}=0.3125\). Their ratio is eight although the equal-torque cross gains are reciprocal. Different normalization denominators explain the apparent directional asymmetry. Ratios are dimensionless only for commensurate coordinates, and rescaling coordinates changes them. Mixed rotational and translational coordinates require explicit units or a justified metric. The more-than-twelvefold quiet-standing index difference reported in the abstract of \citep{Challis2011} should therefore remain attached to that study's definition and model; it does not establish a golf or neural-control strategy.

\subsection{Coordinates and Physical Responses}
Under a constant invertible change \(q=Tz\), virtual work and kinetic energy give
\[
 M_z=T^{\mathsf T}M_qT,\quad Q_z=T^{\mathsf T}Q_q,\quad J_z=J_qT.
\]
Consequently \(J_zM_z^{-1}Q_z=J_qM_q^{-1}Q_q\). A consistently transformed force increment has the same acceleration at the same physical point. The analogous constrained result follows by transforming \(A_z=A_qT\). Representation dependence is a reason to report coordinates and units, not a suggestion that physical predictions are arbitrary.

A nonlinear change needs acceleration transport. For a one-dimensional free particle described by \(q=z^2\), with \(z>0\),
\[
 \ddot q=2z\ddot z+2\dot z^2=0,
 \qquad \ddot z=-\dot z^2/z.
\]
A velocity-dependent coordinate term appears without a physical force or particle acceleration. Complete physical outputs remain consistent when all transport terms are retained. This is one reason to avoid equating the size of a coordinate velocity term with a universal amount of mechanical assistance.

\section{Instantaneous and Accumulated Effects}
\label{sec:instantaneous_cumulative}

\subsection{Integrating a Ledger Along One Trajectory}
If a complete acceleration partition is evaluated along one nominal motion, integration gives
\[
 \dot q(t)-\dot q(t_0)=\sum_k\int_{t_0}^{t}a_k(s)\,ds
                         +\int_{t_0}^{t}a_c(s)\,ds.
\]
This is valid bookkeeping for coordinate-velocity change on smooth intervals, with any impulse jumps accounted for separately. The terms can be signed, oppose one another, or exceed the net change in magnitude. A percentage obtained by dividing by a nearly zero net acceleration is unstable and can be misleading.

The velocity-dependent term is an \emph{instantaneous function of the current state}. That state reflects earlier forces and initial conditions, but it does not uniquely identify their history. For the double integrator \(\ddot q=u\), starting from rest at \(q=0\), both \(u_1(t)=1\) and \(u_2(t)=2+6t^2-6t\) reach \((q,\dot q)=(1/2,1)\) at \(t=1\). Both the input integral and its first time moment agree, despite different inputs. A single state cannot distinguish them.

\subsection{Why Removing an Input Is a Different Calculation}
For a transparent nonlinear example, let \(\dot x=x^2+u\), \(x(0)=0\), and \(u=1\). Before \(t=\pi/2\), \(x(t)=\tan t\). The nominal input integral is \(t\) and the nominal drift integral is \(\tan t-t\); they sum exactly. Removing the input yields \(x(t)=0\), so the intervention difference is \(\tan t\), not the input ledger entry \(t\). At \(t=0.5\), these are approximately 0.5463 and 0.5.

The removed input changes subsequent states and hence subsequent drift. In a swing, it may also change muscle states, support reactions and contact transitions. A force-removal simulation must specify those rules and integrate the resulting system. Neither an instantaneous term nor its nominal integral supplies that simulation automatically. Describing drift as the accumulated effect of all past torques obscures this distinction.

\section{What Throwing Studies Establish}
\label{sec:throwing_lessons}
\citep{Hirashima2008kinetic} studied six baseball players with a 13-degree-of-freedom model comprising trunk, upper arm, forearm, and hand with ball. It partitioned reconstructed generalized-force dynamics into joint-torque, velocity-dependent and gravity terms. It reported larger joint-torque terms at proximal coordinates and substantial velocity-dependent terms in elbow extension and wrist flexion. These were estimates for net generalized moments, not separately measured muscle forces.

The study also integrated nominal induced accelerations and checked reconstructed motion. Its elbow-torque reversal and distal acceleration terms illustrate coupling; they do not identify a protective neural intention. The paper explicitly discusses omitted segments and lower-body/ground interactions. Marker acquisition at 200 Hz and spline resampling at 2000 Hz are different quantities: interpolation does not supply ten times the independent measurement bandwidth.

These findings motivate a golf study but do not settle it. A two-handed club, compliant shaft and foot contacts change the mechanical system. Golf comparisons need matched outputs, force conventions and contact treatment; a reported pitching joint-angle term cannot stand in for clubhead tangential acceleration. The chapter supplies no numerical human-golf finding that late-swing velocity terms universally dominate.

\section{Walking, Muscle Function and Energy}
\label{sec:walking_lessons}

\subsection{Acceleration and Energy Answer Different Questions}
In the 1.5 m/s walking simulation of \citep{Neptune2001}, both soleus and gastrocnemius supported and advanced the trunk during parts of stance, yet their modeled energy delivery differed: soleus delivered energy to the trunk, while gastrocnemius delivered most of its generated energy to the leg during the late-stance/pre-swing interval. Calling gastrocnemius exclusively a swing initiator would hide its simultaneous support and progression roles. The study used both acceleration and segment energetics to make these distinctions.

An acceleration vector alone does not determine energy transfer. Generalized-force power is \(Q_k^{\mathsf T}v\). At a body interface, use the wrench and that body's twist at the same point and in compatible frames: \(F\cdot v_P+\tau_P\cdot\omega\). An ideal fixed support can redirect momentum while doing zero work at its stationary contact point. In a multibody system, individual segments can exchange energy through internal forces even when the net internal work on the whole system cancels. Track the chosen boundary and velocities explicitly.

\subsection{Clinical Examples Are Conditional Model Results}
The sit-to-stand study of \citep{Caruthers2016} used static optimization to estimate muscle forces before induced-acceleration analysis. Its reported forward center-of-mass terms for gluteus maximus and soleus differed from the opposing quadriceps term. That is a useful example of how anatomy alone does not fix the sign of a whole-body output response. It remains conditional on the muscle estimation and model; the present chapter does not independently reproduce it.

IAA has also been applied to stiff-legged gait \citep{RileyKerrigan1999}. A decomposition can help formulate patient-specific hypotheses, but a computed term does not by itself identify a treatable impairment or predict the effect of rehabilitation. Such claims require clinical evidence beyond algebraic closure. The transferable lesson for golf is to test the force-estimation and contact assumptions, not to import a walking or clinical result as a swing prescription.

\section{A Worked Golf-Model Investigation}
\label{sec:golf_implications}

\subsection{A Fully Specified Two-Link Calculation}
Use the uniform rods defined above, with gravity \(g=9.81\) m/s\(^2\), no damping, no other external loads and no additional constraints. This is a toy arm--club analogue, not anthropometry. Angles remain \emph{relative joint coordinates}; do not substitute formulas for two absolute rod angles. The right-hand-side terms are
\[
 \begin{aligned}
 Q_v&=\begin{bmatrix}
 \tfrac12\sin q_2(2v_1v_2+v_2^2)\\-\tfrac12\sin q_2v_1^2
 \end{bmatrix},\\
 Q_g&=-g\begin{bmatrix}
 \tfrac32\sin q_1+\tfrac12\sin(q_1+q_2)\\
 \tfrac12\sin(q_1+q_2)
 \end{bmatrix}.
 \end{aligned}
\]
At \(q=(-45,-90)\) degrees, \(v=(500,200)\) degrees/s and \(Q_u=(10,-2)\) Nm, first convert angles and rates to radians. The joint acceleration contributions in rad/s\(^2\) are
\[
 \begin{aligned}
 a_u&=(9.000,-15.000)^{\mathsf T},\\
 a_v&=(-55.973,170.205)^{\mathsf T},\\
 a_g&=(7.804,2.601)^{\mathsf T},\\
 \ddot q&=(-39.170,157.806)^{\mathsf T}.
 \end{aligned}
\]
The distal endpoint is \(p=(\sin q_1+\sin(q_1+q_2),-\cos q_1-\cos(q_1+q_2))^{\mathsf T}\) m. Differentiating this map gives the following point-acceleration terms in m/s\(^2\):
\[
 \begin{aligned}
 Ja_u&=(10.607,-2.121)^{\mathsf T},\\
 Ja_v&=(-120.353,-41.195)^{\mathsf T},\\
 Ja_g&=(-1.839,-12.876)^{\mathsf T},\\
 \dot Jv&=(159.394,-51.695)^{\mathsf T},\\
 \ddot p&=(47.808,-107.887)^{\mathsf T}.
 \end{aligned}
\]
The large positive horizontal curvature term and negative velocity-force term partly cancel. Omitting either changes the answer qualitatively. The gravitational point term is not simply free-fall acceleration because the proximal pin constrains the linked system. These numbers demonstrate why the output projection and complete accounting matter more than an isolated label such as proximal, distal or passive.

\subsection{From the Toy Model to Human Evidence}
The prior golf study of \citep{Takagi2021pelvis}, published online in 2019 and in a 2021 journal issue, used 31 skilled golfers, motion capture and two force plates to partition pelvis angular acceleration. Its existence rules out a broad first-golf-application claim here. Its abstract supports this study description; this chapter does not claim to reproduce its full methods or extend its pelvis findings to the clubhead.

A defensible next investigation would combine measured kinematics and support loads with a specified golfer--club model. Report the output point and frame, net-moment or muscle-force estimation, grip and foot contact assumptions, residuals and solver details. Then calculate complete acceleration terms and vary plausible inertial, kinematic, actuator and contact parameters. Sensitivity to differentiation noise is especially relevant because inverse dynamics and output acceleration use time derivatives of measured motion.

Use closure as a necessary implementation check: the sum must match a separate complete forward solve within a stated absolute and relative tolerance. Closure alone is not validation against humans; even an inaccurate model can close algebraically. Compare predicted motion or held-out observables with independent measurements and propagate uncertainty to signed output terms. Keep uncertain terms visible rather than assigning the residual to a favored muscle or mechanism.

\section{Superposition and Its Limits}
\label{sec:superposition_limits}
Three distinct results should remain distinguishable. An \emph{instantaneous decomposition} explains one calculated acceleration under a stated partition. A \emph{local sensitivity} reports the change in an output for a small feasible input change at that state. A \emph{trajectory intervention} predicts a changed motion after integrating a specified counterfactual. The first does not supply the third; the second remains local to its contact mode.

There is also a distinction between acceleration authority and energetic cost. Two postures can have different \(JPB\) maps for the same generalized input while requiring different muscle forces, activation, stability margins and contact reactions. A favorable gain alone therefore does not prove reduced human effort, safety or superior technique. Those are additional outputs with additional evidence requirements.

\section{Summary}
\label{sec:iaa_summary}
Induced acceleration analysis links a complete force model to a declared motion output. Its core calculation is a coupled solve at one state. Active constraints supply an affine baseline, and point acceleration adds kinematic curvature. Posture changes the coupling map; consistent coordinate transformations preserve physical output increments. Nominal integrals account for a velocity change but do not reconstruct a unique force history or replace a force-removal simulation. Muscle roles require force-estimation evidence, and energy flow requires power and boundary accounting. Together these distinctions make IAA a useful bridge between biomechanics and control theory.



\begin{exercises}{Chapter Exercises}{ex:iaa}
\begin{enumerate}
\item Reproduce the two unit-torque responses at \(q_2=0\) and the zero cross-response at \(q_2=\arccos(-2/3)\). Compute both normalized ratios and explain why their asymmetry does not contradict the symmetric inverse mass matrix.
\item For the circular-guide example, apply \(F=(4,6)\) N. Verify \(\ddot q=(-9,3)\) m/s\(^2\) and \(\lambda=-22\) N. Split the applied force into its two Cartesian components and show that adding two complete constrained solutions would double the radial baseline. Explain why reaction feasibility must be checked before interpreting a unilateral version of this example.
\item Reproduce the worked two-link acceleration ledger and differentiate its endpoint map to recover all five point vectors. Compute the endpoint velocity and project each point term onto its unit tangent. State why the resulting signed speed-rate terms differ from vector magnitudes.
\item Integrate both double-integrator input histories from \(t=0\) to \(t=1\). Verify the common final state using \(\dot q(1)=\int_0^1u(t)dt\) and \(q(1)=\int_0^1(1-t)u(t)dt\). Explain which force-history information that state fails to preserve.
\item Design a comparison between a net-joint-moment golf model and a muscle-driven model. Specify a shared physical output, contact mode, residual convention and independent measurement. Identify which conclusions need acceleration, power, muscle-force estimation or an integrated intervention.
\end{enumerate}
\end{exercises}
\flushbottom{}
